% Wahrscheinlichkeit – bank/wahrscheinlichkeit/ (zusammenbau v0.4, Heft)
\einheitenkopf[t2]{Wahrscheinlichkeit}
\setcounter{aufgabe}{5}

% Anlauf (ohne Prüfkennung)
\begin{aufgabe}{Zählen – Anlauf}
\begin{teile}
\teil Paul hat zwei Mützen (grau, grün) und zwei Schals (gelb, weiß). Schreibe alle Paare aus Mütze und Schal auf – erst alle mit der grauen Mütze.
\teil Lena hat 3 Hosen und 2 Pullover. Wie viele Kombinationen aus Hose und Pullover gibt es? \leerfeld Kombinationen
\teil Eine Münze wird geworfen (Kopf oder Zahl) und ein Glücksrad mit den Feldern 1, 2, 3 gedreht. Schreibe alle Ergebnisse als Paare (Münze; Glücksrad) auf.
\end{teile}
\end{aufgabe}

\begin{aufgabe}{Zählen – Prüfungsaufgaben}
\begin{teile}
\teil Die Ziffern 1, 4 und 7 werden in zufälliger Reihenfolge zu einer dreistelligen Zahl gelegt. Schreibe alle möglichen Zahlen auf. Welcher Anteil davon ist gerade? \hfill (P10 2016 OS)
\teil Auf drei Kärtchen stehen die Ziffern 5, 1 und 7. Welche ist die größte dreistellige Zahl, die man damit legen kann? \hfill (P10 2016 OS)
\teil Zwei gleiche Spielsteine mit einer roten und einer blauen Seite werden gleichzeitig geworfen. Wie viele verschiedene Ergebnisse sind möglich? \hfill (P10 2017 OS)
\teil Zwei Glücksräder haben je vier gleich große Felder: links 4, 5, 4, 6, rechts 5, 5, 6, 4. Beide werden gedreht; gelesen wird erst die linke, dann die rechte Ziffer. Schreibe alle möglichen zweistelligen Zahlen auf. \hfill (P10 2025 OS)
\teil Ein Würfel wird zweimal geworfen. Schreibe alle Paare auf, bei denen der zweite Wurf eine 4 zeigt. \hfill (P10 2020 OS)
\teil Von den fünf Punkten P, Q, R, S, T liegen Q, R und S auf einer Geraden; sonst liegen keine drei auf einer Geraden. Wie viele verschiedene Dreiecke mit drei dieser Punkte als Ecken gibt es? Kreuze an. \hfill (P10 2015 OS) \\ \kreuz{7} \\ \kreuz{8} \\ \kreuz{9} \\ \kreuz{11}
\steil Ein Fahrradschloss hat einen dreistelligen Code aus den Ziffern 2, 4 und 8; jede Ziffer kommt genau einmal vor. Wie viele Codes sind möglich? \hfill (P10 2017 OS)
\steil Lina hat den Code ihres Spinds vergessen. Er besteht aus den Ziffern 1, 3 und 6, jede genau einmal. Wie viele Codes muss sie höchstens ausprobieren? \hfill (P10 2017 OS)
\end{teile}
\end{aufgabe}

% Anlauf (ohne Prüfkennung)
\begin{aufgabe}{Einstufig – Anlauf}
\begin{teile}
\teil In einer Lostrommel liegen 30 Nieten und 5 Gewinne. Gesucht ist die Wahrscheinlichkeit für einen Gewinn. Unterstreiche, was möglich und was günstig ist; rechne nicht. möglich: \leerfeld , günstig: \leerfeld
\teil Ein Würfel wird einmal geworfen. Wie groß ist die Wahrscheinlichkeit für eine 4?
\teil Das Glücksrad hat acht gleich große Felder. Wie groß ist die Wahrscheinlichkeit für Blau?

\kreisdiagramm{rot/1,blau/1,rot/1,gelb/1,rot/1,blau/1,gelb/1,rot/1}
\end{teile}
\end{aufgabe}

\begin{aufgabe}{Einstufig – Prüfungsaufgaben}
\begin{teile}
\teil In einer Schachtel liegen 50 Murmeln: 18 blaue, 9 rote und 23 grüne. Eine Murmel wird gezogen. Wie groß ist die Wahrscheinlichkeit für Rot als Bruch, als Dezimalzahl und in Prozent? \hfill (P10 2019 OS)
\teil Ein Glücksrad hat 8 gleich große Felder: 2-mal Stern, 5-mal Mond, 1-mal Sonne. Wie groß ist die Wahrscheinlichkeit für Mond als Bruch und in Prozent? \hfill (P10 2014 OS)
\teil Ein Glücksrad hat acht gleich große Felder mit den Zahlen 0 bis 7. Wie groß ist die Wahrscheinlichkeit für die 0 als Bruch und in Prozent? \hfill (P10 2017 OS)
\teil Bei einer Tombola liegen 45 Nieten und 15 Gewinnlose in der Trommel. Wie groß ist die Wahrscheinlichkeit für einen Gewinn? \hfill (P10 2014 OS)
\teil Lose tragen die Nummern 251 bis 600; genau ein Los gewinnt ein Fahrrad. Ole sagt: „Mit einem Los ist die Wahrscheinlichkeit für das Fahrrad $\frac{1}{350}$.“ Hat er Recht? Begründe. \hfill (P10 2016 OS)
\teil Eine Kiste enthält 60 Eier, 4 davon sind zerbrochen. Ein Ei wird zufällig genommen. Wie groß ist die Wahrscheinlichkeit, dass es zerbrochen ist? \hfill (P10 2016 OS)
\teil Ein Glücksrad hat acht gleich große Felder mit den Zahlen 1 bis 8. Wie groß ist die Wahrscheinlichkeit für weder 1 noch 8? \hfill (P10 2014 OS)
\teil Ein Glücksrad hat 4 gleich große Felder in Rot und Gelb. Die Wahrscheinlichkeit für Gelb ist 75\,\%. Wie viele Felder sind gelb? \hfill (P10 2018 OS)
\end{teile}
\end{aufgabe}

\begin{aufgabe}{Einstufig – Prüfungsaufgaben – weiter}
\begin{teile}
\teil Ein Glücksrad hat gleich große Felder; 6 davon sind rot, die übrigen weiß. Die Wahrscheinlichkeit für Rot soll $\frac{3}{4}$ sein. Zeichne das ganze Glücksrad mit den fehlenden weißen Feldern. \hfill (P10 2019 OS)

\kreisleer[2]
\teil Auf einem Teller liegen 18 Muffins, 3 davon mit Chili, die anderen mit Schokolade. Jonas hat schon 2 Schoko-Muffins gegessen. Mara sagt: „Die Wahrscheinlichkeit, jetzt einen Chili-Muffin zu erwischen, ist $\frac{3}{18}$.“ Hat sie Recht? Begründe. \hfill (P10 2018 OS)
\steil Bei einem Schulfest gibt es Lose mit den Nummern 201 bis 700. Einen Trostpreis gibt es für jedes Los, dessen Nummer auf 45 endet, außer für das Los 445 (Hauptgewinn). Timo hat ein Los und sieht nur die letzte Ziffer: 5. Er sagt: „Die Wahrscheinlichkeit für einen Trostpreis ist $\frac{4}{50}$.“ Hat er Recht? Begründe. \hfill (P10 2016 OS)
\steil Bei einer Tombola gibt es Lose mit den Nummern 201 bis 800. Einen Trostpreis gibt es für jedes Los, dessen Nummer auf 72 endet, außer für das Los 572 (Hauptgewinn). Sina hat ein Los und sieht nur die letzte Ziffer: 2. Sie sagt: „Die Wahrscheinlichkeit für einen Trostpreis ist $\frac{6}{60}$.“ Hat sie Recht? Begründe. \hfill (P10 2016 OS)
\end{teile}
\end{aufgabe}

% Anlauf (ohne Prüfkennung)
\begin{aufgabe}{Topf mit P = 50 \% auswählen – Anlauf}
\begin{teile}
\teil Glücksrad Stern hat 6 gleich große Felder, 3 davon blau; Glücksrad Mond hat 9 gleich große Felder, 3 davon blau; Glücksrad Sonne hat 8 gleich große Felder, 3 davon blau. Bei welchem Glücksrad ist die Wahrscheinlichkeit für Blau $\frac{1}{3}$? Kreuze an. \\ \kreuz{Stern} \\ \kreuz{Mond} \\ \kreuz{Sonne}
\end{teile}
\end{aufgabe}

\begin{aufgabe}{Topf mit P = 50 \% auswählen – Prüfungsaufgaben}
\begin{teile}
\teil Drei Dosen enthalten je 8 Bonbons: links 5 Kirsch und 3 Zitrone, in der Mitte 4 Kirsch und 4 Zitrone, rechts 2 Kirsch und 6 Zitrone. Aus welcher Dose zieht man Kirsch mit der Wahrscheinlichkeit 50\,\%? Kreuze an. \hfill (P10 2015 OS) \\ \kreuz{linke Dose} \\ \kreuz{mittlere Dose} \\ \kreuz{rechte Dose}
\end{teile}
\end{aufgabe}

% Anlauf (ohne Prüfkennung)
\begin{aufgabe}{Mit Zurücklegen – Anlauf}
\begin{teile}
\teil Aus einer Dose wird ein Keks genommen und gegessen, dann ein zweiter. Mit oder ohne Zurücklegen? Kreuze an. \\ \kreuz{mit Zurücklegen} \\ \kreuz{ohne Zurücklegen}
\teil Eine Münze wird zweimal geworfen. Trage im Baum an jeden Ast die Wahrscheinlichkeit ein.

\baumzwei{Kopf/,Zahl/}{Kopf/,Zahl/}{Kopf/,Zahl/}
\teil Ein Basketballspieler trifft beim Freiwurf mit der Wahrscheinlichkeit $\frac{4}{5}$; er wirft zweimal. Ergänze die fehlenden Wahrscheinlichkeiten im Baum.

\baumzwei{Treffer/\frac{4}{5},daneben/}{Treffer/,daneben/\frac{1}{5}}{Treffer/\frac{4}{5},daneben/}
\end{teile}
\end{aufgabe}

\begin{aufgabe}{Mit Zurücklegen – Prüfungsaufgaben}
\begin{teile}
\teil Ein grüner Würfel trägt die Zahlen 2, 3, 3, 5, 6, 6, ein gelber Würfel 1, 2, 4, 4, 6, 6. Erst wird der grüne, dann der gelbe geworfen; betrachtet wird gerade oder ungerade. Ergänze die fehlenden Wahrscheinlichkeiten im Baum. \hfill (P10 2026 FOR)

\baumzwei{gerade/\frac{3}{6},ungerade/}{gerade/\frac{5}{6},ungerade/}{gerade/,ungerade/}
\teil Ein Glücksrad mit 5 gleich großen Feldern (3-mal grün, 2-mal gelb) wird zweimal gedreht. Im Baum ist nur die Wahrscheinlichkeit für Gelb auf der ersten Stufe eingetragen. Ergänze alle fehlenden Wahrscheinlichkeiten. \hfill (P10 2014 OS)

\baumzwei{grün/,gelb/\frac{2}{5}}{grün/,gelb/}{grün/,gelb/}
\teil In einer WG liegen jede Woche drei Zettel im Topf: Bad, Küche, frei; alle Zettel kommen jede Woche zurück. Lara zieht jede Woche als Erste. Sie sagt: „Die Wahrscheinlichkeit, zwei Wochen hintereinander ‚frei‘ zu ziehen, ist $\frac{1}{9}$.“ Zeige, dass das stimmt. \hfill (P10 2024 OS)
\steil Zwei Glücksräder haben je vier gleich große Felder: links 7, 8, 7, 7, rechts 8, 9, 8, 7. Beide werden gedreht; gelesen wird erst links, dann rechts. Wie groß ist die Wahrscheinlichkeit für 78 oder 87? \hfill (P10 2025 OS)
\steil Ein Glücksrad hat 8 gleich große Felder: 5-mal rot, 2-mal blau, 1-mal grün. Es wird zweimal gedreht. Wie groß ist die Wahrscheinlichkeit, zweimal dieselbe Farbe zu drehen? \hfill (P10 2014 OS)
\teil Ein roter Würfel trägt die Zahlen 1, 2, 2, 3, 5, 6, ein blauer Würfel 2, 3, 3, 5, 5, 5. Erst wird der rote, dann der blaue geworfen. Wie groß ist die Wahrscheinlichkeit, dass nicht zweimal eine ungerade Zahl fällt? \hfill (P10 2026 FOR)
\end{teile}
\end{aufgabe}

\begin{aufgabe}{Mit Zurücklegen – Prüfungsaufgaben – weiter}
\begin{teile}
\teil Zwei Glücksräder haben je vier gleich große Felder: links 5, 6, 5, 7, rechts 6, 7, 6, 5. Beide werden gedreht; gelesen wird erst links, dann rechts. Eva behauptet: „66 ist genauso wahrscheinlich wie 77.“ Hat Eva Recht? Begründe. \hfill (P10 2025 OS)
\teil Ein Würfel wird zweimal geworfen. Mia behauptet: „Zwei gleiche Augenzahlen sind wahrscheinlicher als eine 1 im ersten Wurf.“ Hat Mia Recht? Begründe mit Wahrscheinlichkeiten. \hfill (P10 2020 OS)
\steil Zwei leere Glücksräder haben je vier gleich große Felder; gelesen wird erst das linke, dann das rechte. Trage Ziffern so ein, dass die Wahrscheinlichkeit für 57 genau $\frac{3}{8}$ ist, und zeige es mit einer Rechnung. \hfill (P10 2025 OS)

\kreisdiagramm{/1,/1,/1,/1} \kreisdiagramm{/1,/1,/1,/1}
\teil Ein roter Würfel trägt die Zahlen 1, 1, 3, 4, 5, 6, ein blauer Würfel 2, 3, 3, 4, 5, 6. Erst wird der rote, dann der blaue geworfen. Ändere genau eine Zahl auf dem blauen Würfel so, dass die Wahrscheinlichkeit für die Augensumme 3 genau $\frac{1}{9}$ ist, und zeige es mit einer Rechnung. \hfill (P10 2026 FOR)
\steil Jana dreht ein Glücksrad mit 4 gleich großen Feldern (1 Stern, 3 leer) dreimal. Sie gewinnt, wenn mindestens zweimal der Stern kommt. Beschrifte alle Äste des Baums und berechne die Gewinnwahrscheinlichkeit. \hfill (P10 2020 OS)

\baumdreigleich{Stern/,leer/}
\steil Ole wirft einen Würfel dreimal. Er gewinnt, wenn mindestens zweimal eine Zahl größer als 4 fällt. Beschrifte alle Äste des Baums und berechne die Gewinnwahrscheinlichkeit. \hfill (P10 2020 OS)

\baumdreigleich{größer als 4/,höchstens 4/}
\end{teile}
\end{aufgabe}

% Anlauf (ohne Prüfkennung)
\begin{aufgabe}{Ohne Zurücklegen – Anlauf}
\begin{teile}
\teil In einer Dose sind 9 Kekse, 4 davon mit Schoko. Ein Schokokeks wird gegessen. Wie viele Kekse sind noch da, wie viele davon mit Schoko? \leerfeld insgesamt, davon \leerfeld
\teil In einer Urne liegen 3 rote und 5 blaue Kugeln. Zwei Kugeln werden ohne Zurücklegen gezogen. Wie groß ist die Wahrscheinlichkeit für zweimal Rot?
\teil In einem Beutel liegen 4 rote und 3 grüne Kugeln. Zwei werden ohne Zurücklegen gezogen. Ergänze die fehlenden Wahrscheinlichkeiten im Baum.

\baumzwei{rot/\frac{4}{7},grün/}{rot/,grün/\frac{3}{6}}{rot/,grün/}
\end{teile}
\end{aufgabe}

\begin{aufgabe}{Ohne Zurücklegen – Prüfungsaufgaben}
\begin{teile}
\steil In einem Hut liegen 7 Zettel für die Aufgaben einer Woche, auf einem steht „frei“. Tom zieht zuerst, dann zieht Lea aus den übrigen Zetteln. Wie groß ist die Wahrscheinlichkeit, dass Lea „frei“ zieht? \hfill (P10 2024 OS)
\steil In einer Schachtel liegen 15 Pralinen, 5 davon mit Nuss. Emil nimmt nacheinander drei Pralinen, ohne zurückzulegen. Ergänze die beiden leeren Felder im Baum. \hfill (P10 2019 OS)

\baumdrei{Nuss/\frac{5}{15},ohne/}{Nuss/\frac{4}{14},ohne/\frac{10}{14}}{Nuss/\frac{5}{14},ohne/\frac{9}{14}}{Nuss/\frac{3}{13},ohne/\frac{10}{13}}{Nuss/\frac{4}{13},ohne/\frac{9}{13}}{Nuss/,ohne/\frac{9}{13}}{Nuss/\frac{5}{13},ohne/\frac{8}{13}}
\steil Beim Staffellauf wird die Reihenfolge der 9 Läuferinnen zufällig ausgelost; jede läuft genau einmal. Wie groß ist die Wahrscheinlichkeit, dass Pia unter den ersten drei ist? \hfill (P10 2015 OS)
\steil In einer Kiste liegen 15 Stifte: 5 rote, 3 grüne und 7 schwarze. Drei Stifte werden ohne Zurücklegen gezogen. Leo behauptet: „Drei rote sind doppelt so wahrscheinlich wie drei grüne.“ Widerlege die Behauptung mit einer Rechnung. \hfill (P10 2019 OS)
\steil Auf einem Blech liegen 18 Berliner: 15 mit Pflaume, 3 mit Vanille. Ida nimmt ohne hinzusehen nacheinander zwei. Schreibe alle Ergebnisse auf, bei denen mindestens ein Vanille-Berliner dabei ist. Berechne die Wahrscheinlichkeit, dass beide mit Pflaume gefüllt sind. Welches Ereignis gehört zur Rechnung $\frac{3}{18} \cdot \frac{2}{17} + \frac{15}{18} \cdot \frac{14}{17}$? \hfill (P10 2018 OS)
\steil In einer Schale liegen 15 Pralinen: 13 mit Nougat, 2 mit Marzipan. Tom nimmt ohne hinzusehen nacheinander zwei. Schreibe alle Ergebnisse auf, bei denen mindestens eine Marzipan-Praline dabei ist. Berechne die Wahrscheinlichkeit, dass beide mit Nougat gefüllt sind. Welches Ereignis gehört zur Rechnung $\frac{2}{15} \cdot \frac{13}{14} + \frac{13}{15} \cdot \frac{2}{14}$? \hfill (P10 2018 OS)
\end{teile}
\end{aufgabe}
